\documentclass[10pt,twocolumn,a4paper]{article}

\usepackage[utf8]{inputenc}
\usepackage[T1]{fontenc}
\usepackage{times}
\usepackage{graphicx}
\usepackage{amsmath}
\usepackage{amssymb}
\usepackage{geometry}
\usepackage{url}

\geometry{
 a4paper,
 left=1.6cm,
 right=1.6cm,
 top=1.8cm,
 bottom=2.0cm,
 columnsep=0.7cm
}

\setlength{\parindent}{0.35cm}
\setlength{\parskip}{0pt}

\title{\textbf{Project Title: A Concise and Informative Scientific Title}}

\author{
 Student Name Surname 1, Student Name Surname 2\\
 Student Name Surname 3, Student Name Surname 4\\
 Bachelor's Degree Programme in Applied Computer Science and Artificial Intelligence (ACSAI)\\
 Course: AI-LAB -- Computer Vision, Signal Processing, and Natural Language Processing\\
 Department of Computer Science, Sapienza University of Rome\\
 Matricola: XXXXXXX, XXXXXXX, XXXXXXX, XXXXXXX\\
 \texttt{<student1@studenti.uniroma1.it>}\\
 \texttt{<student2@studenti.uniroma1.it>}\\
 \texttt{<student3@studenti.uniroma1.it>}\\
 \texttt{<student4@studenti.uniroma1.it>}
}

\date{}

\begin{document}
 
 \maketitle
 
 \begin{abstract}
  This report presents the project developed for the AI-LAB course: Computer Vision, Signal Processing, and Natural Language Processing. The abstract must briefly summarize the problem, the motivation, the proposed approach, the dataset, the experimental protocol, and the main results. It should also clearly state the main contribution of the project, for example a model adaptation, a new pipeline, a meaningful tuning strategy, an experimental comparison, or a critical analysis of an existing method. The abstract should be self-contained and should not include unnecessary implementation details.
 \end{abstract}
 
 \noindent\textbf{Keywords:} artificial intelligence; deep learning; computer vision; signal processing; natural language processing.
 
 \section{Introduction}
 
 The introduction should motivate the problem and explain why it is relevant in the context of Artificial Intelligence. It should not be a generic description of the topic, but a clear explanation of the specific problem addressed by the project.
 
 Students should clearly discuss:
 \begin{itemize}
  \item the problem addressed by the project;
  \item why the problem is relevant;
  \item the main limitations of existing approaches;
  \item what the project tries to improve, modify, test, compare, or analyze;
  \item the personal contribution of the students with respect to existing papers, repositories, tutorials, or baseline methods.
 \end{itemize}
 
 The contribution does not necessarily have to be a completely new model. It may consist of a careful benchmark, a meaningful modification of an existing method, a new preprocessing or post-processing step, a different training strategy, an ablation study, a robustness analysis, or a critical comparison between methods.
 
 The last paragraph of this section should briefly summarize the structure of the report.
 
 \section{Related Work}
 
 This section should present the most relevant previous works. Students should avoid a simple list of papers. Each cited work should be connected to the project and discussed in relation to the proposed approach.
 
 A good related work section should include:
 \begin{itemize}
  \item classical or early approaches related to the problem;
  \item recent deep learning or AI-based methods;
  \item relevant datasets, benchmarks, or evaluation protocols;
  \item limitations of the existing literature;
  \item how the proposed project differs from previous work.
 \end{itemize}
 
 All external resources must be cited properly, including papers, datasets, GitHub repositories, tutorials, pretrained models, and libraries when they are central to the project. For example, students may cite foundational deep learning work~\cite{lecun2015deep}, convolutional architectures~\cite{he2016resnet}, transformer-based models~\cite{vaswani2017attention}, or task-specific references when appropriate.
 
 \section{Methodology}
 
 This section should describe the proposed method or pipeline. The description must be precise enough to allow another student or researcher to reproduce the work.
 
 The first element of this section should be a figure spanning both columns, showing the complete architecture or pipeline. The figure should include the input, preprocessing, main model blocks, output, and evaluation flow.
 
\begin{figure*}[t]
 \centering
 \includegraphics[width=0.95\textwidth]{01-ArchitectureExample.png}
 \caption{Example of an architecture and experimental pipeline. The figure spans both columns and illustrates the expected level of clarity for the methodology section: input data, preprocessing, feature extraction or backbone, proposed contribution, prediction head, training loop, and evaluation protocol. Students should replace this example with a diagram of their own method.}
 \label{fig:architecture}
\end{figure*}
 
 Students should describe:
 \begin{itemize}
  \item the overall architecture or pipeline shown in Fig.~\ref{fig:architecture};
  \item the input data and preprocessing operations;
  \item the model or models used;
  \item what is taken from existing literature or implementations;
  \item what has been modified, tuned, added, or analyzed by the students;
  \item the loss function or objective function;
  \item the training strategy and implementation details.
 \end{itemize}
 
 If the project is based on an existing architecture or repository, this must be stated explicitly. The report must clearly separate what is standard, what is reused, and what is the students' own contribution.
 
 \section{Experimental Setup}
 
 This section describes the experimental protocol. It should be concise, reproducible, and directly connected to the research question.
 
 \subsection{Dataset}
 
 This subsection should describe the dataset used in the project. Students should report only the information that is useful to understand the experiments.
 
 Important information includes:
 \begin{itemize}
  \item dataset name and source;
  \item number of samples and classes;
  \item input format;
  \item train, validation, and test split;
  \item preprocessing steps;
  \item relevant limitations or biases;
  \item whether the dataset was used as provided or modified by the students.
 \end{itemize}
 
 \begin{table}[h]
  \centering
  \caption{Dataset summary.}
  \label{tab:dataset}
  \begin{tabular}{|l|c|}
   \hline
   \textbf{Property} & \textbf{Value} \\
   \hline
   Dataset name & XXXX \\
   Number of samples & XXXX \\
   Number of classes & XX \\
   Training samples & XXXX \\
   Validation samples & XXXX \\
   Test samples & XXXX \\
   Input format & XXXX \\
   \hline
  \end{tabular}
 \end{table}
 
 \subsection{Training and Evaluation Protocol}
 
 This subsection should describe how the experiments were performed.
 
 Students should specify:
 \begin{itemize}
  \item hardware and software environment;
  \item main libraries and frameworks;
  \item optimizer, learning rate, batch size, and number of epochs;
  \item loss function;
  \item evaluation metrics;
  \item baseline methods;
  \item hyperparameter tuning strategy, if any;
  \item any procedure used to ensure a fair comparison.
 \end{itemize}
 
 \begin{table}[h]
  \centering
  \caption{Training configuration.}
  \label{tab:training}
  \begin{tabular}{|l|c|}
   \hline
   \textbf{Parameter} & \textbf{Value} \\
   \hline
   Framework & XXXX \\
   Optimizer & XXXX \\
   Learning rate & XXXX \\
   Batch size & XXXX \\
   Epochs & XXXX \\
   Loss function & XXXX \\
   Main metric & XXXX \\
   \hline
  \end{tabular}
 \end{table}
 
 \subsection{Results}
 
 This subsection should present the quantitative and qualitative results. Students should avoid reporting numbers without explanation.
 
 Depending on the task, possible metrics include:
 \begin{itemize}
  \item accuracy, precision, recall, and F1-score for classification;
  \item mAP for object detection;
  \item Dice score or IoU for segmentation;
  \item MAE, RMSE, or correlation for regression;
  \item inference time, memory usage, model size, or FLOPs when efficiency is relevant.
 \end{itemize}
 
 \begin{table}[h]
  \centering
  \caption{Comparison between baseline and proposed approach.}
  \label{tab:results}
  \begin{tabular}{|l|c|c|c|}
   \hline
   \textbf{Method} & \textbf{Metric 1} & \textbf{Metric 2} & \textbf{Time} \\
   \hline
   Baseline & XX & XX & XX \\
   Proposed method & XX & XX & XX \\
   \hline
  \end{tabular}
 \end{table}
 
 When useful, students should include qualitative examples, confusion matrices, error cases, visual predictions, or comparisons between different configurations.
 
 \subsection{Critical Discussion}
 
 This subsection is essential. Students must interpret the results and explain what they mean.
 
 The discussion should answer:
 \begin{itemize}
  \item which method performs best and why;
  \item whether the proposed modification or analysis produced a real benefit;
  \item which classes, samples, or conditions are more difficult;
  \item where the model fails;
  \item whether the results are coherent with the related work;
  \item what the main limitations of the project are;
  \item what could be improved in future work.
 \end{itemize}
 
 The discussion should connect the numerical results to the contribution stated in the introduction. A strong report is not only based on high performance, but on a clear and critical interpretation of the experiments.
 
 \section{Conclusion}
 
 The conclusion should summarize the project without repeating the whole report.
 
 It should briefly restate:
 \begin{itemize}
  \item the problem addressed;
  \item the method or pipeline used;
  \item the main experimental findings;
  \item the students' personal contribution;
  \item the main limitations and possible future developments.
 \end{itemize}
 
 \section*{Acknowledgments}
 
 Optional. Students may acknowledge collaborators, datasets, libraries, tools, or external resources used in the project.
 
 \begin{thebibliography}{9}
  
  \bibitem{lecun2015deep}
  Y. LeCun, Y. Bengio, and G. Hinton,
  ``Deep learning,''
  \textit{Nature}, vol. 521, no. 7553, pp. 436--444, 2015.
  
  \bibitem{he2016resnet}
  K. He, X. Zhang, S. Ren, and J. Sun,
  ``Deep residual learning for image recognition,''
  in \textit{Proceedings of the IEEE Conference on Computer Vision and Pattern Recognition},
  pp. 770--778, 2016.
  
  \bibitem{vaswani2017attention}
  A. Vaswani, N. Shazeer, N. Parmar, J. Uszkoreit, L. Jones, A. N. Gomez, L. Kaiser, and I. Polosukhin,
  ``Attention is all you need,''
  in \textit{Advances in Neural Information Processing Systems},
  2017.
  
  \bibitem{lin2014coco}
  T.-Y. Lin, M. Maire, S. Belongie, J. Hays, P. Perona, D. Ramanan, P. Dollar, and C. L. Zitnick,
  ``Microsoft COCO: Common objects in context,''
  in \textit{European Conference on Computer Vision},
  pp. 740--755, 2014.
  
 \end{thebibliography}
 
\end{document}
